\section{Induction-free representation of certified computations}
\label{sec:p05-representability}

The strong representation needed by the diagonal construction is the
pointwise statement with an unrestricted output quantifier: for each
\emph{standard} input tuple \(\vec a\),
\(\mathsf Q\derives\forall y(\gamma_f(\overline{\vec a},y)
\leftrightarrow y=\overline{f(\vec a)})\).
Proving all correct and incorrect \emph{numeral} instances alone would not
exclude a nonstandard output. Conversely, this pointwise assertion is not
the universal termination assertion
\(\mathsf Q\derives\forall\vec x\exists!y\gamma_f(\vec x,y)\).
The construction below proves the former without asserting the latter.
All complexity assertions in this section use the handbook's declared
\emph{marked} arithmetic presentation from
\cref{def:language-arithmetic} and the following syntactic-hierarchy
definition: marked \(\leq\)-quantifiers are counted as bounded before
their eliminable comparison-witness expansion. In particular, a literal
existential prefix with bounded matrix here means that displayed marked
presentation; it is not a claim about the quantifier prefix of the fully
expanded raw \(\Lang_A\) string.
For a comparison of these quantifier distinctions and the use of finite
comparison in \(\mathsf Q\), see \cite[chs.~11,16--17,39]{Smith2020QRepresentation}.
The certificate construction here is a separately specified reconstruction;
it does not identify its formula with the \(\PA\) graphs of
\cref{thm:p05-pa-pr-graphs}.

\begin{hblemma}{thm:p05-rep-bounded-decisions}{proved=annex:p05-rep-bounded-decisions}{Fixed bounded computations in \(\mathsf Q\)}
For any bounded arithmetic formula \(A(\vec x)\) in the marked syntax of
\cref{def:language-arithmetic}, and any standard tuple \(\vec a\),
\(\mathsf Q\) proves \(A(\overline{\vec a})\) when it is true in
\(\mathcal N\), and proves its negation when it is false.
For each fixed standard \(n\), bounded quantification over
\(u\leq\bar n\) is provably equivalent in \(\mathsf Q\) to the
corresponding finite conjunction or disjunction of instances at
\(\bar0,\ldots,\bar n\). The analogous assertion for \(u<\bar n\)
uses \(\bar0,\ldots,\overline{n-1}\), with empty disjunction false and
empty conjunction true.
\end{hblemma}

\begin{hbdefinition}{def:p05-rep-register-certificate}{A fixed program and its bounded trace formula}
Fix a finite primitive recursive description \(d\) of arity \(k\),
using exactly the initial functions, composition and primitive recursion
of \cref{sec:turing-machines-02}. The annex gives an effective compilation
to a finite program \(M_d\) on registers \(R_0,\ldots,R_K\), where
\(K\geq k\). Inputs are in \(R_0,\ldots,R_{k-1}\), the output is in
\(R_k\), and the remaining initial registers are zero.
Its finitely many instruction labels have initial label \(0\).
Instructions are increment-and-jump, zero-test/decrement-and-jump, and halt.
These are numerical programs used to specify certificates, not additions
to the arithmetic language or replacements for the program numbering in
Part IV.

Use the already expanded bounded formulas \(\mathsf E\) and
\(\mathsf E^{(2)}\) of \cref{def:p04-tagged-history-entry}. For this
certificate only, family \(20\) stores the transition count at position
\(0\), family \(21\) the instruction label at time \(i\), and family
\(22\) the register \(j\) at time \(i\). Define
\[
\begin{aligned}
 \operatorname{Row}_d(s,i,a_0,\ldots,a_K,p)
 &\defeq \mathsf E(s,\overline{21},i,p)\land
       \bigwedge_{j=0}^K\mathsf E^{(2)}(s,\overline{22},i,\bar j,a_j),\\
 \operatorname{Le}(u,v)&\defeq\exists w\leq v\,(w+u=v).
\end{aligned}
\]
The finite conjunction is a metalinguistic expansion, not variable-arity
arithmetic syntax. \(\operatorname{Le}(u,v)\) has only a marked bounded
quantifier and logically implies the source abbreviation \(u\leq v\).

The quantifier-free formula
\(\operatorname{Step}_d(p,\vec a,q,\vec b)\) is the finite disjunction
over the nonhalt instructions of \(M_d\). An instruction at label \(h\)
which increments register \(j\) and jumps to \(h'\) contributes
\[
 p=\bar h\land q=\overline{h'}\land b_j=Sa_j
                       \land\bigwedge_{v\ne j}b_v=a_v.
\]
A zero-test/decrement instruction with zero target \(h_0\) and positive
target \(h_1\) contributes
\[
 p=\bar h\land\bigwedge_{v\ne j}b_v=a_v\land
 \bigl((a_j=0\land b_j=0\land q=\overline{h_0})\lor
       (a_j=Sb_j\land q=\overline{h_1})\bigr).
\]
Let \(\operatorname{Halt}_d(p)\) be the finite disjunction of its halt
labels, and \(\operatorname{Label}_d(p)\) the finite disjunction of all
its labels. Use fresh variables in every abbreviation expansion. Put
\[
\begin{aligned}
 C_d(\vec x,y,s)\defeq {}&\operatorname{Le}(y,s)\land
 \exists\ell\leq s\,\bigl[\mathsf E(s,\overline{20},0,\ell)\land
 \exists a_0\leq s\cdots\exists a_K\leq s\,
    (\operatorname{Row}_d(s,0,\vec a,0)\land
       \bigwedge_{j<k}a_j=x_j\land\bigwedge_{k\leq j\leq K}a_j=0)\\
 &\land\forall i\leq\ell\,\exists p\leq s\,
                         (\mathsf E(s,\overline{21},i,p)
                                      \land\operatorname{Label}_d(p))\\
 &\land\forall i\leq\ell\,\bigl(i=\ell\lor
   \exists p\leq s\,\exists q\leq s\,
   \exists a_0\leq s\cdots\exists a_K\leq s\,
   \exists b_0\leq s\cdots\exists b_K\leq s\,
    (\operatorname{Row}_d(s,i,\vec a,p)\land
     \operatorname{Row}_d(s,Si,\vec b,q)\land
                     \operatorname{Step}_d(p,\vec a,q,\vec b))\bigr)\\
 &\land\exists p\leq s\,
       (\mathsf E(s,\overline{21},\ell,p)\land\operatorname{Halt}_d(p)
                         \land\mathsf E^{(2)}(s,\overline{22},\ell,\bar k,y))
 \bigr].
\end{aligned}
\]
Every displayed conjunction indexed by register numbers expands to a fixed
finite conjunction. The other quantifiers are marked bounded with
arithmetic-term bounds; after expanding \(\mathsf E,\mathsf E^{(2)}\),
no array-access, pairing, subtraction or program-execution function symbol
occurs. No standalone comparison is being treated as an atomic bounded
formula.
\end{hbdefinition}

\begin{hbtheorem}{thm:p05-rep-trace}{proved=annex:p05-rep-trace}{Standard adequacy of the bounded trace certificate}
For each fixed description \(d\), \(C_d\) is bounded and
\[
 f_d(\vec a)=b\quad\Longleftrightarrow\quad
       \mathcal N\entails\exists s\,C_d(\overline{\vec a},\bar b,s).
\]
Its numerical truth characteristic, with all arguments standard numbers,
is primitive recursive. Moreover \(\mathsf Q\) proves
\(C_d(\vec x,y,s)\to y\leq s\).
The adequacy assertion is in the standard structure; it does not assert
internal termination or unique decoding for all nonstandard certificates.
\end{hbtheorem}

\begin{hbdefinition}{def:p05-rep-least-graph}{The selected least-certificate graph}
Define the bounded matrix and its literal existential-prefix graph by
\[
\begin{aligned}
 \delta_d(\vec x,y,s)&\defeq C_d(\vec x,y,s)\land
       \forall t\leq s\,
         \bigl(t=s\lor\neg\exists z\leq t\,C_d(\vec x,z,t)\bigr),\\
 \gamma_d(\vec x,y)&\defeq\exists s\,\delta_d(\vec x,y,s).
\end{aligned}
\]
Thus \(\gamma_d\) is literally an existential quantifier applied to a
bounded formula in that declared marked presentation, not merely
standard-model equivalent to a marked prefix formula.
The minimum is taken over certificate numbers, not over computation times
or output values. The marked guard \(\forall t\leq s\,H(t)\) uses exactly
the source abbreviation
\(\forall t((\exists u\,(u+t=s))\to H(t))\), with fresh \(u\).
It does \emph{not} replace its comparison by
\(\operatorname{Le}(t,s)\). Consequently a premise
\(\bar r\leq s\) obtained from the fixed-standard dichotomy supplies
the guard's instance at \(\bar r\) directly by universal elimination
and implication elimination. No bounded-witness conversion is required.
The logical expansion just given exposes the meaning used in \(\mathsf Q\)
derivations; it does not reclassify its comparison witnesses as marked
bounded quantifiers. The complexity assertion instead follows directly
from the original marked guard and the hierarchy's generating clauses.
\end{hbdefinition}

\begin{hbtheorem}{thm:p05-rep-strong-graph}{proved=annex:p05-rep-strong-graph}{Pointwise strong graph representation in \(\mathsf Q\)}
For every fixed primitive recursive description \(d\) and every standard
tuple \(\vec a\),
\[
 \mathsf Q\derives\forall y\,
   \bigl(\gamma_d(\overline{\vec a},y)
                         \leftrightarrow y=\overline{f_d(\vec a)}\bigr).
\]
The formula also defines the actual numerical graph in \(\mathcal N\).
In particular the universally quantified output is not restricted to
numerals. No universal input-totality conclusion for this formula is made.
\end{hbtheorem}

\begin{hbcorollary}{thm:p05-rep-numeralwise-relations}{proved=annex:p05-rep-numeralwise-relations}{Decidable numerical relations and finite aggregates}
If \(R\subseteq\N^k\) has a fixed primitive recursive characteristic
description \(d_R\), with value \(1\) for truth and \(0\) for falsehood,
then
\(\rho_R(\vec x)\defeq\gamma_{d_R}(\vec x,\bar1)\)
is literally \(\exists s\,\delta_{d_R}(\vec x,\bar1,s)\), has a bounded
matrix, is accurate in \(\mathcal N\), and numeralwise represents \(R\)
in \(\mathsf Q\), positively and negatively. For every fixed standard
bound, any finite list of these instance proofs can be assembled into the
corresponding bounded universal statement. This last assertion does not
require \(\rho_R\) itself to be bounded.
\end{hbcorollary}

\begin{hbdefinition}{def:p05-rep-certified-presentation}{Selection for the governing certified proof codes}
Retain the syntax numbering and derivation checks of
\cref{thm:arithmetization-syntax}. In particular, a certified proof has
the governing code \(N_{16}(N_{15}(R),W)\): \(R\) is its checked
postorder derivation record and \(W\) has exactly one accepting history
for each \emph{root open axiom occurrence}, in the order computed by the
checker. Fix the Part IV index \(e_T\) whose domain is exactly the set
of codes of the given c.e. axioms of \(T\). Each history is checked by
the bounded \(\mathsf{Comp}(\overline{e_T},a,s)\) from
\cref{thm:bounded-verification}. Temporary discharged assumptions do not
need axiom histories; an empty open list requires \(W=0\).
This fixes the c.e. certificate convention, even if the chosen theory
happens to have a finite axiom set.

Let \(d_T\) be the finite primitive recursive description of the
characteristic of this \emph{actual numerical} certified proof checker,
including comparison with the conclusion code \(x\). Select
\[
 \mathsf{Proof}^{*}_T(p,x)\defeq\gamma_{d_T}(p,x,\bar1)
                 =\exists s\,\kappa_T(p,x,s),\qquad
 \kappa_T(p,x,s)\defeq\delta_{d_T}(p,x,\bar1,s),
\]
and \(\mathsf{Prov}^{*}_T(x)\defeq\exists p\exists s\,
\kappa_T(p,x,s)\). The star denotes an explicit selection available in
\cref{def:fixed-proof-predicate}, not a theorem that every previously
chosen representative is \(\mathsf Q\)-equivalent to it. Proofs using
this selection must use it throughout.
\end{hbdefinition}

\begin{hbtheorem}{thm:p05-rep-certified-proof}{proved=annex:p05-rep-certified-proof}{The c.e. certified-proof interface}
Let \(T\) be a c.e. axiom set whose theory extends \(\mathsf Q\), with
the fixed index and certificate convention just specified. The numerical
proof relation is primitive recursive, and \(\mathsf{Proof}^{*}_T\)
is literally existential-prefix with bounded matrix, is accurate in
\(\mathcal N\), and has both positive and negative \(\mathsf Q\) proofs
on every standard numeral pair. For every sentence \(\sigma\),
\[
 T\derives\sigma\quad\Longleftrightarrow\quad
 \mathcal N\entails\exists p\,\mathsf{Proof}^{*}_T
                       (p,\overline{\ulcorner\sigma\urcorner}).
\]
For each fixed standard \(n,a\), a finite list of negative proof-instance
decisions gives
\(\mathsf Q\derives\forall p\leq\bar n\,
\neg\mathsf{Proof}^{*}_T(p,\bar a)\), whenever every numerical
candidate \(0,\ldots,n\) is rejected. Invalid syntactic, derivation,
conclusion, or certificate codes are among these rejected inputs.
No consistency or soundness assumption is needed for these coding facts.
\end{hbtheorem}

\begin{hbdefinition}{def:p05-rep-diagonal-default}{Total diagonal substitution with a fixed default}
Fix distinguished variable \(x_0\). On a formula code \(n\) with free
variables contained in \(\{x_0\}\), let \(\mathsf{diag}(n)\) be the
code of numeral substitution \(x_0:=\bar n\). Put
\(\tau=\ulcorner0=0\urcorner\). On every other number set
\(\mathsf{diag}(n)=\tau\). This specifies the previously arbitrary default
without changing any intended diagonal instance. Its computation uses the
governing numeral function \(\nu(0)=N_1()\),
\(\nu(n+1)=N_2(\nu(n))\), and the capture-avoiding substitution clauses
of \cref{thm:arithmetization-syntax}. Also set
\(\mathsf{neg}(n)=N_7(n)\) on formula codes and \(N_7(\tau)\) otherwise.
Neither numerical function is a new primitive arithmetic symbol.
\end{hbdefinition}

\begin{hbcorollary}{thm:p05-rep-diagonal-graph}{proved=annex:p05-rep-diagonal-graph}{The diagonal and negation graph interface}
Both \(\mathsf{diag}\) and \(\mathsf{neg}\) have fixed primitive recursive
descriptions. For their selected graphs, in particular
\(D(x,y)=\gamma_{d_{\rm diag}}(x,y)=\exists s\,
\delta_{d_{\rm diag}}(x,y,s)\) with bounded matrix, every standard
\(n\), including invalid codes, satisfies
\[
 \mathsf Q\derives\forall y\,
      (D(\bar n,y)\leftrightarrow
                              y=\overline{\mathsf{diag}(n)}).
\]
Thus existence and unrestricted-output uniqueness at the fixed numeral
input used by the diagonal lemma are supplied together. The analogous
assertion holds for negation coding.
\end{hbcorollary}

\begin{hbwarning}{def:p05-rep-boundary}{Which facts this interface does not supply}
Standard-model adequacy of the finite certificate is not a universal
\(\mathsf Q\) trace-determinism or sequence-existence theorem. The
least-certificate guard is what rules out nonstandard competing outputs
at a fixed standard input. The selected \(\mathsf{Proof}^{*}_T\) is
\(\Sigma^0_1\), not claimed bounded; the embedded axiom-history check
\(\mathsf{Comp}\) is bounded, which is a different statement.
Numeralwise decisions and fixed finite aggregates suffice for the usual
Rosser comparison. They do not make its formula automatically
\(\Pi^0_1\). Finally, the \(\PA\) derivability conditions and internal
proof transformations require their separate universal invariant proofs;
the present pointwise interface is not a replacement for them.
\end{hbwarning}
